% Mathematical fragment for deterministic clause-derived witness discharge.
% Prepared in this session by the logical-reduction subagent.
% Inherited sources: incoming ipra_witness_discharge_2026-11-11,
% especially report.tex sections 2.1--2.8 and validation/macros.py.
% The new result is the deterministic two-premise rule, not the scalar bridge.

\section{Discharging a witness using its first clause instance}
\label{sec:deterministic-name}

The incoming witness-discharge presentation has a separate proof of an
equality $b=t$ inside each eliminated witness scope. That premise is
unnecessary. It suffices to recognize an equality in one fixed instance of
the clause which the witness itself supplies. If that instance has no usable
equality, the witness is discharged without substitution. The choice is
determined by the major formula; the certificate supplies neither an
equality proof nor a replacement term.

This is a restriction of the previous inference, with a nontrivial
completeness theorem. It does not rely on the numerical standard model,
semantic truth testing, or an unlisted equality rule.

\subsection{Complete native specification}

Use one numerical sort. Let $N$ denote full computation normalization of
the scalar terms
\[
 t ::= x\mid 0\mid St\mid I(n,a;u.h)
\]
for the context-closed reductions
\[
 I(0,a;u.h)\longrightarrow a,
 \qquad I(Sn,a;u.h)\longrightarrow h[I(n,a;u.h)/u].
 \tag{NF}
\]
Only the state variable $u$ in $h$ is bound. Every displayed term means
its normal form, and every substitution is capture-avoiding and followed
by $N$. Normalization of a formula acts on numerical terms and leaves its
formula constructors intact. We use the established termination,
confluence, and substitution-coherence results for these reductions.
The normalization algorithm is part of the trusted syntax and must be
included in a complete-checker measurement.

Formulas are
\[
 A ::= (s=t)\mid D_{x,y}A,
\]
where $x,y$ are distinct bound numerical variables. A reusable context
clause is $[\vec A;y.B]$, where $\vec A$ has zero or one antecedent,
and $y$ is absent from every antecedent. Its ordinary reading is
\[
 [\,;y.B]^\circ=\forall y B^\circ,
 \qquad [A;y.B]^\circ=A^\circ\to\forall y B^\circ.
\]
An ordinary assumption $A$ abbreviates $[\,;v.A]$, with $v$ unused.
Contexts are finite lists of clause occurrences. Every proof starts with
the single explicit initial clause
\[
                 \mathcal R=[\,;y.y=y]. \tag{Ax}
\]
Put $H=(0=0)$, $1=S0$, and $2=S1$. The designated contradiction is $0=1$.
There are exactly the following three primitive inference families.

\paragraph{Use.}
For a retained $[\vec A;y.B]$ and any well-scoped term $t$, infer $N(B[t/y])$
after proving the antecedent if there is one. There is no proof premise
for a clause with no antecedent.
\[
 \frac{\{\Gamma\vdash A:A\text{ is the sole antecedent, if present}\}}
      {\Gamma\vdash N(B[t/y])}.
 \tag{Use}
\]

\paragraph{Intro.}
For fresh $z$ and a well-scoped witness $t$,
\[
 \frac{\Gamma,N(B(t,\eta))\vdash N(B(t,Sz))}
      {\Gamma\vdash D_{x,y}B(x,y)},
 \qquad
 \eta=\begin{cases}
 z,& B\text{ is an equation and }t\text{ is literally }0,\\
 0,&\text{otherwise}.
 \end{cases}
 \tag{Intro}
\]
Here $z$ is fresh for $\Gamma,t$ and all fixed parameters of $B$.
The displayed ordinary assumption is discharged.

\paragraph{Name.}
After proving $D_{x,y}B$, take a fresh numerical parameter $b$ and form
\[
 B_b(y)=N(B(b,y)),\qquad
 A_0=N(B_b(0)),\qquad A_1=N(B_b(1)),
 \qquad \mathcal B_b=[A_0;y.B_b(y)].
\]
Define the replacement term \emph{deterministically} by
\[
 \tau_b(B)=
 \begin{cases}
 t,& A_0\text{ is literally }H\text{ and }A_1
       \text{ is literally the equation }b=t,\\
 b,&\text{otherwise}.
 \end{cases}
 \tag{First-instance test}\label{eq:first-instance-test}
\]
The literal equation test is ordered: a formula $t=b$ is not recognized
unless its left side is itself literally $b$. The term $t$ is the right
side of $A_1$, so no separate term argument is supplied. Then the rule is
\[
 \frac{\Gamma\vdash D_{x,y}B(x,y)\qquad
       \Gamma,\mathcal B_b\vdash F}
      {\Gamma\vdash N(F[\tau_b(B)/b])},
 \tag{Name}\label{eq:deterministic-name}
\]
provided the exported normal form has no free $b$. The parameter $b$ is
fresh for $\Gamma$ and every fixed parameter of $B$. The branch formula
$F$ is well scoped in the extended numerical context. A recognized $t$
may contain $b$; substitution is simultaneous, so an inserted copy of $t$
is not substituted again.

These are all the primitive inferences. Reflexivity is a Use of
$\mathcal R$. Computation is accounted for by canonical syntax. There
are no additional primitive equality, logical, induction, case, cut, or
weakening rules.

\subsection{Soundness and the escape condition}

The ordinary formula reading is
\[
 (s=t)^\circ=(s=t),\qquad
 (D_{x,y}B)^\circ=
       \exists x\bigl(B^\circ(x,0)\to\forall y B^\circ(x,y)\bigr).
\]
For this subsection, use intuitionistic first-order equality logic with
the scalar computation equations, numerical coverage, separation
$0=1\to\bot$, and contextual induction for literal equations.
This is the inherited scalar theory with a full-iPRA bi-interpretation.

\begin{lemma}[The replacement equality is always available]
\label{lem:det-replacement}
In the context $\Gamma,\mathcal B_b$, one can derive
$b=\tau_b(B)$ using only the initial reflexivity clause and Use.
\end{lemma}
\begin{proof}
If the test fails, $\tau_b(B)=b$ and Use of $\mathcal R$ at $b$ proves
$b=b$. If it succeeds, the antecedent of $\mathcal B_b$ is exactly
$H$. Use of $\mathcal R$ at zero proves $H$, and Use of $\mathcal B_b$
at $1$ proves its normalized first instance, which is exactly
$b=\tau_b(B)$. No equality reasoning occurs in either derivation.
\end{proof}

\begin{proposition}[Soundness of Name]
Rule~\eqref{eq:deterministic-name} is sound for the ordinary reading,
including when the recognized replacement term contains $b$.
\end{proposition}
\begin{proof}
Ordinary existential elimination opens the major formula at $b$ with
the assumption represented by $\mathcal B_b$. The branch gives $F$.
Lemma~\ref{lem:det-replacement} gives $b=\tau_b(B)$, and ordinary equality
substitution gives $F[\tau_b(B)/b]$. Computation equations identify this
with the exported normal form. That normal form contains no free $b$,
so ordinary existential elimination discharges the witness and its clause.
Equality substitution itself does not require $b$ to be absent from its
right side; only existential discharge needs the stipulated absence from
the final exported formula.
\end{proof}

Use has its displayed ordinary reading. Intro is sound by numerical
coverage in case mode and literal-equation induction in progress mode.
Thus the entire presentation is sound in the inherited scalar theory.

\subsection{Ordinary elimination and numerical cases without equality}

\begin{lemma}[Ordinary paired elimination]
\label{lem:det-ordinary-elim}
The usual rule
\[
 \frac{\Gamma\vdash D_{x,y}B\qquad
       \Gamma,\mathcal B_b\vdash A}
      {\Gamma\vdash A},\qquad b\notin\operatorname{FV}(A),
\]
is derivable with one use of Name.
\end{lemma}
\begin{proof}
Apply Name with its supplied branch formula $A$. Since $A$ is a normal
formula with no free $b$, substituting any term for $b$ leaves it unchanged.
Thus $N(A[\tau_b(B)/b])=A$ regardless of whether the first-instance test
succeeds. There is no need to force the test to fail or to choose a target.
\end{proof}

Weakening, injective fresh renaming, and ordinary-assumption cut are
admissible by induction on proof trees. Freshen local eigenparameters
before inserting proofs. The deterministic test depends only on the
major formula and is unchanged by an injective renaming. An ordinary
assumption is an unconditional constant clause, so its Use nodes can be
replaced by the supplied proof. None of these structural constructions
uses general equality transport or arbitrary numerical substitution of
proof trees.

For unused $u$, write $q_y A(y)=D_{u,y}A(y)$.
\begin{lemma}[Conditional closure and numerical cases]
\label{lem:det-cases}
The calculus derives
\[
 \frac{\Gamma,A(0)\vdash A(Sz)}{\Gamma\vdash q_yA(y)},
 \qquad
 \frac{\Gamma\vdash q_yA(y)\qquad\Gamma\vdash A(0)}
      {\Gamma\vdash A(t)},
\]
and consequently
\[
 \frac{\Gamma\vdash A(0)\qquad\Gamma\vdash A(Sz)}
      {\Gamma\vdash A(t)}
 \tag{Cases}\label{eq:det-cases}
\]
for every native formula $A$ and appropriate fresh $z$.
These constructions use no equality transport or equation induction.
\end{lemma}
\begin{proof}
For the first rule, introduce the unused witness at the numeral $1$.
It is not literally zero, so Intro uses the case premise at every matrix
root. For the second, use Lemma~\ref{lem:det-ordinary-elim}. The opened
clause is $[A(0);y.A(y)]$, because the witness coordinate is unused.
Supply its antecedent and apply Use at $t$. The result has no free
existential witness, so ordinary elimination closes. For Cases, first
weaken the successor proof by the unused ordinary assumption $A(0)$,
then combine the preceding two constructions.
\end{proof}

\subsection{Recovering equality from the fixed first instance}

Define normalizing selector and parity terms by
\[
 C(v,a,c)=N(I(v,a;u.c)),\qquad u\notin\operatorname{FV}(c),
 \qquad P(y)=N(I(y,0;u.C(u,1,0))).
\]
Direct computation gives literal identities of canonical syntax:
\[
 C(0,a,c)\equiv a,\qquad C(Sv,a,c)\equiv c,
 \qquad P(1)\equiv1,\qquad P(2)\equiv0.
 \tag{Endpoints}\label{eq:det-endpoints}
\]
Only these fixed endpoints of $P$ are used; no parity induction or range
theorem is required.

\begin{theorem}[Uncrossed term transport]
\label{thm:det-transport}
For every fixed numerical template $h(w)$, the rule
\[
 \frac{\Gamma\vdash s=t\qquad\Gamma\vdash N(h(s))=r}
      {\Gamma\vdash N(h(t))=r}
\]
is derivable. It uses no new arithmetic assumption and no supplied
equality branch at Name.
\end{theorem}
\begin{proof}
Choose $x,y$ fresh for all passive data. Let $h(x)$ denote its canonical
instance and form the equation matrix
\[
 \begin{split}
 B(x,y)=\bigl(&C(y,0,C(P(y),N(h(x)),x))\\
             &{}=C(y,0,C(P(y),r,t))\bigr).
 \end{split}
 \tag{Naming matrix}\label{eq:det-naming-matrix}
\]
Introduce $D_{x,y}B(x,y)$ at witness $s$. Its successor instance is
\[
 C(P(Sz),N(h(s)),s)=C(P(Sz),r,t).
\]
Apply Cases to the template
\[
       A(v)=\bigl(C(v,N(h(s)),s)=C(v,r,t)\bigr)
\]
at target $P(Sz)$. Its zero instance is exactly the supplied equation
$N(h(s))=r$, and its successor instance is exactly the supplied equation
$s=t$. The eigenparameter used for this Cases argument is fresh for the
current local context and is different from $z$.

The resulting successor proof ignores the ordinary assumption of Intro.
Consequently it is valid whether $s$ is literally zero, selecting progress,
or any other term, selecting cases. This constructs the major formula
without any equality substitution.

Open that major at fresh $b$. By the endpoint computations,
\[
 B(b,0)\equiv H,\qquad B(b,1)\equiv(b=t),
 \qquad B(b,2)\equiv(N(h(b))=r).
\]
The first-instance test therefore selects exactly $\tau_b(B)=t$.
For the sole branch of Name, prove $H$ by the fixed reflexivity clause
and specialize the opened clause at $2$. This yields
$F=(N(h(b))=r)$. Name exports
\[
 N\bigl((N(h(b))=r)[t/b]\bigr)
       \equiv (N(h(t))=r)
\]
by normalization--substitution coherence. The parameter $b$ was fresh
for $h,t,r$, so the exported formula is outer scoped.

The premise carrying $s=t$ is used while proving the major naming formula,
not as a third premise of Name. Thus the proof does not smuggle the deleted
equality branch back into the rule's checking interface.
\end{proof}

\begin{corollary}[All equality reasoning]
\label{cor:det-equality}
Symmetry, transitivity, congruence, crossed transport, and transport in
every native formula template are derivable.
\end{corollary}
\begin{proof}
For symmetry, take $h(w)=w$ and $r=s$ in
Theorem~\ref{thm:det-transport}; the minor premise is reflexivity $s=s$.
This proves $t=s$ from $s=t$. Uncrossed transport followed by symmetry
gives crossed transport. A reflexive minor gives congruence, followed by
symmetry if necessary. Symmetry of $a=b$ and transport of $b=c$ through
the identity template give $a=c$.

For transport in an equation $p(w)=q(w)$, transport its left side,
reverse the resulting equation, transport its right side, and reverse
again. This proves $p(t)=q(t)$ from $s=t$ and $p(s)=q(s)$.

For a compound template $D_{x,y}E(w,x,y)$, argue by external induction
on formula depth. Open $D_{x,y}E(s,x,y)$ at a fresh $c$ using ordinary
elimination. Introduce $D_{x,y}E(t,x,y)$ at witness $c$. Since $c$ is
a fresh parameter, its syntax is not literally zero, and Intro selects
cases; no proposition $c\ne0$ is assumed. From the temporary base
$E(t,c,0)$, use the induction hypothesis and symmetry to obtain
$E(s,c,0)$. Use the opened old clause at $Sz$ to obtain $E(s,c,Sz)$,
then transport forward by the induction hypothesis to $E(t,c,Sz)$.
Close Intro and ordinary elimination. Every recursively transported
matrix has smaller formula depth. This is not induction over quantified
formulas inside the arithmetic theory.
\end{proof}

\subsection{Exact equivalence with unrestricted witness discharge}

Let $\mathcal W$ be the incoming canonical witness-discharge calculus:
the same terms, formulas, clauses, fixed axiom and Use/Intro rules, with
\[
 \frac{\Gamma\vdash D_{x,y}B\qquad
       \Gamma,\mathcal B_b\vdash b=t\qquad
       \Gamma,\mathcal B_b\vdash F}
      {\Gamma\vdash N(F[t/b])}
 \tag{Old Name}
\]
and the same freshness and export conditions. Let $\mathcal A$ denote
the deterministic calculus just specified.

\begin{theorem}[Identity-language sequent equivalence]
\label{thm:det-equivalence}
For every finite canonical clause context $\Gamma$ containing the
distinguished reflexivity axiom and every well-scoped canonical formula
$A$,
\[
               \Gamma\vdash_{\mathcal A} A
          \quad\Longleftrightarrow\quad
               \Gamma\vdash_{\mathcal W} A.
\]
The two translations are effective transformations of finite derivations
and preserve the complete end sequent literally.
\end{theorem}
\begin{proof}
From $\mathcal A$ to $\mathcal W$, translate the two premises of each
new Name recursively and supply the missing equality proof from
Lemma~\ref{lem:det-replacement}. All other rules are shared.

In the converse direction, recursively translate the three premises of
an Old Name inference. In its extended witness context, use
Corollary~\ref{cor:det-equality} to transport the translated branch
formula along the translated equation $b=t$. This gives the prescribed
normal form $N(F[t/b])$. Its original export condition says that it has
no free $b$, so Lemma~\ref{lem:det-ordinary-elim} closes the major
existential. This works even when $t$ contains $b$, because the same
ordinary equality substitution and final export condition apply.

The apparent dependency is not circular: ordinary elimination and
Cases were obtained before any equality transport; the naming-matrix
proof of term transport uses only those constructions, the supplied
premises, and one new Name. Formula transport then uses strictly
smaller matrix depth. Every replacement is therefore a finite, explicitly
defined proof construction.
\end{proof}

\begin{corollary}[Full iPRA, not just the equational fragment]
The ordinary realization of $\mathcal A$ is bi-interpretable and
equiconsistent with full first-order intuitionistic primitive recursive
arithmetic. Both numerical interpretation domains are whole, equality
is actual equality, and both numerical composite maps have graph $x=y$.
\end{corollary}
\begin{proof}
Compose Theorem~\ref{thm:det-equivalence} with the incoming canonical
calculus equivalence and its inherited scalar arithmetic and logical
bi-interpretation. The present translation is the identity on the entire
native language and every end sequent, so it preserves both inherited
formula round trips and the numerical identity maps. In particular it
preserves the designated contradiction $0=1$, yielding equiconsistency.
The scalar compiler, equality reflection, predecessor, logical macros,
and arithmetic bridge are inherited and must be included or precisely
referenced; the local result does not replace their proofs.
\end{proof}

\subsection{Implementation invariants and claim boundaries}

A forward checker may infer the major formula and the branch formula.
At Name it checks that the major has a $D$ root, opens its matrix at the
newest fresh parameter, computes its zero and first instances, and
determines $\tau_b(B)$ by~\eqref{eq:first-instance-test}. It then checks
the branch under the one new clause, performs simultaneous free-parameter
substitution, normalizes, and verifies that the result has no free $b$.
The original outer goal is compared with this inferred result.

Scope validation is essential even if repeated formula-grammar checks are
removed: a branch may use the fresh witness in an equation, but Name may
export only an outer-scoped formula. The clause antecedent must be the
actual normalized zero instance of the same opened matrix; replacing it
by a claimed true formula would be unsound. The recognized first instance
must be an equation with left side literally the current witness. Merely
having an equation root, being provably equal to that shape, or having the
witness on the right does not authorize the substitution in this rule.

There is a useful implementation invariant: if the computed zero instance
is literally $H$, then the matrix has an equation root, and so does its
first instance. Normalization never deletes a $D$ constructor. A checker
which has already validated the major formula may therefore omit a second
equation-root test on the first instance; the test $A_0=H$ already ensures
that its two equation fields exist. This is a structural fact about the
validated grammar, not a relaxation of the mathematical guard.

Numerical substitution can change the selected Intro mode or create the
first-instance guard. A primitive-tag-preserving substitution theorem is
not claimed. Admissibility follows through the exact equivalence with the
reference calculus and its established substitution theorem or repairs.

The simplification deletes one proof premise and two certificate payloads
(the replacement term and, for a synthesizing checker, the claimed branch
formula) from the earlier interface. It still has three inference families
and one explicit reflexivity clause. A measured checker reduction must
count the deterministic guard and normalization in full. It is neither a
proof of absolute minimum size nor a claim that three families first
occurred here. The new mathematical assertion is the exact completeness
of the first-instance restriction.

Uniform full normalization over arbitrary finite iterator descriptions is
total recursive but need not be primitive recursive. A finite budget or
explicit reduction certificates can give a primitive-recursive bounded
certificate predicate without changing the provable formulas; that
interface is distinct and its code must also be counted. No object-level
uniform evaluator or new totality axiom is asserted by this metasyntactic
operation.
